\section{Callbacks Compilation}\label{sec:callbacks}

A distinctive feature of qstack's compiler framework is its principled approach to compiling programs with classical callbacks. In qstack, callbacks are treated as opaque components of the quantum program: the compiler does not analyze or modify their internal logic. Instead, compiler passes wrap existing callbacks in adapter functions that handle --- and compile --- any dynamically generated kernels.

% \subsection{The Callback Wrapping Pattern}

When compiling a kernel, for example to introduce error correction, there are two issues in dealing with the arbitrary, unknown-to-the-compiler callback that the kernel has prior to compilation:
\begin{enumerate}
\item The original callback expects measurement arguments at a certain level of abstraction. If we are implementing a lower level of abstraction, we will need to essentially \emph{decode} the post-compilation measurement results to pass to the original callback
the value(s) it expects.
\item The original callback will optionally return a new kernel and that kernel will need to be compiled.
\end{enumerate}
Both issues are handled by having the compiler create a callback that ``wraps'' the original callback. A decoding computation prior to calling the original callback handles the first issue. A call back to the compiler itself handles the second issue.

Here is how the idiom looks for our majority-vote example where we assume the original callback is called \texttt{original\_callback} and the compiler itself is bound to \texttt{rep3\_compiler}. The new callback, produced by the compiler, could be:
\begin{minted}{python}
def decoder():
  m1 = context.pop()
  m2 = context.pop()
  m3 = context.pop()
  if (m1+m2+m3) >= 2:
    return 1 
  else
    return 0
  
def wrapped_callback(measurement):
  context.push(measurement)
  k = original_callback(decoder())
  if k is None
    return None
  return self.compile(k)
\end{minted}

Notice that once you define a decoding procedure, as done here in \texttt{decoder}, the rest of \texttt{wrapped\_callback} is boilerplate; it encodes the same idiom each compiler phase uses. It can be done manually in a compiler phase, as shown here, but we also provide some straightforward support in our framework implementation for it.
%In practice, the decoding logic need not be ``inlined'' into the new callback but can be a call to a Python function provided by the compiler. The boilerplate of the recursive call to the compiler can also be supported by our framework more directly, but the idea is the same.

The essential elegance is that this callback wrapping composes across compilation phases with each phase using the underlying phase's callbacks by passing them decoded information and then compiling any kernels the underlying phase generates. Thus, at runtime, each generated kernel is compiled by the full sequence of compilation phases without any phase needing to know anything about the sequence.

%The central challenge in compiling quantum programs with callbacks is handling kernels that are generated dynamically at runtime. %Since these kernels do not exist during initial compilation, they must be compiled just-in-time before execution. A compiler pass %must therefore not only transform the initial quantum program but also ensure that any kernels generated by callbacks are %similarly transformed.


%If implemented manually, a compiler pass would need to wrap each callback in an adapter function that calls the original callback %with the measurement result, returns None if the callback returns None, and compiles and returns the transformed kernel if the %callback returns a kernel.

%For more specialized pipelines like those in error correction, in which a decoder must calculate a logical measurement based on %the current context, the compiler also needs a \`{decoder} function that can replace the original physical measurement with the %logical one.

%Here's how this pattern would look if implemented manually:

%\begin{minted}{python}
%# Original callback
%def original_callback(measurement):
%    if measurement == 1:
%        return Kernel("q", [Mix("q")])  # Returns a kernel with toy instructions
%    return None

%# Manual adapter function for toy-to-standard compilation
%def wrapped_callback(measurement):
%    # Call the original callback
%    result = original_callback(decoder(measurement))
    
%    # If no kernel was returned, just return None
%    if result is None:
%        return None
    
%    # Otherwise, compile the returned kernel
%    return toy_compiler.compile_kernel(result)
%\end{minted}

%The key insight is that this adapter doesn't need to understand the internal logic of the callback as it treats the callback as an opaque function that might return a kernel. This wrapping pattern is inherently composable: multiple compiler passes can each wrap callbacks in their own adapters, creating a chain of transformations that are applied to any dynamically generated kernels. For example, error correction codes can be easily concatenated:

%\begin{minted}{python}
%error_corrected_program = \
%  SteaneCompiler().compile(
%    TrivialRepetitionCompiler().compile(original_program)
%  )
%\end{minted}

%When executed, these wrapped callbacks ensure that dynamically generated kernels go through the complete compilation pipeline, with transformations applied in the correct reverse order, the outermost wrapper (from the last compiler in the chain) runs first.